\chapter{Flatulence (Gas)}

\section*{Definition}
Flatulence is the expulsion of intestinal gas through the rectum. It is a normal physiological process, with most adults passing gas 10-25 times daily. Excessive or malodorous flatulence may indicate dietary imbalance, malabsorption, or gastrointestinal pathology.

\section*{What Happens in Your Body}
Intestinal gas is produced by two mechanisms: (1) swallowed air (aerophagia) and (2) bacterial fermentation of undigested carbohydrates in the colon. Fermentation by colonic bacteria produces hydrogen, methane, and carbon dioxide. Sulfur-containing gases (hydrogen sulfide, methanethiol) from protein fermentation produce characteristic odor. Gas accumulates and is expelled when intraluminal pressure exceeds sphincter tone.

\section*{Causes}
\begin{itemize}
  \item Dietary factors: Beans, lentils, cruciferous vegetables, onions, garlic, whole grains, carbonated beverages, chewing gum
  \item Lactose intolerance (deficiency of lactase enzyme)
  \item Fructose malabsorption (impaired fructose transport)
  \item Celiac disease (gluten-induced villous atrophy)
  \item Small intestinal bacterial overgrowth (SIBO)
  \item Irritable bowel syndrome (IBS -- altered gas handling)
  \item Pancreatic insufficiency (reduced digestive enzyme secretion)
  \item Aerophagia (excessive air swallowing from anxiety, rapid eating, gum chewing)
  \item Constipation (prolonged transit increases fermentation time)
\end{itemize}

\section*{When to See a Doctor}
See your doctor if:
\begin{itemize}
  \item Persistent, severe, or worsening flatulence
  \item Associated abdominal pain, bloating, or distension
  \item Diarrhea, steatorrhea, or unintentional weight loss (suggests malabsorption)
  \item Blood in stool or iron deficiency anemia
  \item Sudden change in bowel habits
  \item Infants or elderly with excessive gas and distress
\end{itemize}

\section*{Related Symptoms}
\begin{itemize}
  \item Abdominal bloating
  \item Abdominal pain or cramping
  \item Belching
  \item Diarrhea
  \item Constipation
  \item Borborygmi (stomach growling)
  \item Nausea
\end{itemize}

\section*{Self-Care / Home Management}
\begin{itemize}
  \item Avoid gas-producing foods: beans, cabbage, onions, broccoli, cauliflower, carbonated beverages.
  \item Eat slowly, chew thoroughly, and avoid talking while eating to reduce aerophagia.
  \item Try lactase supplements if lactose intolerant.
  \item Over-the-counter simethicone may help reduce gas bubble formation.
  \item Activated charcoal pads in underwear can reduce odor.
  \item Keep a food-symptom diary to identify trigger foods.
\end{itemize}

\section*{How Your Doctor Will Evaluate You}
\begin{itemize}
  \item History and dietary review to identify triggers.
  \item Lactose hydrogen breath test for lactose intolerance.
  \item Fructose hydrogen breath test for fructose malabsorption.
  \item Glucose hydrogen breath test for SIBO.
  \item Stool studies (fecal elastase for pancreatic insufficiency, stool culture).
  \item Celiac serology (tissue transglutaminase IgA).
  \item Colonoscopy if alarm features present (bleeding, weight loss, anemia).
\end{itemize}

\section*{Treatment Options}
\begin{itemize}
  \item Dietary modification: Low-FODMAP diet under dietitian guidance.
  \item Digestive enzyme supplements (lactase, alpha-galactosidase, pancreatic enzymes).
  \item Antibiotics for SIBO (rifaximin).
  \item Probiotics may help rebalance gut flora.
  \item Treatment of underlying condition (gluten-free diet for celiac, pancreatic enzyme replacement for insufficiency).
\end{itemize}


\section*{References}
\begin{thebibliography}{99}
\bibitem{levitt1980}
Levitt MD. Volume and composition of human intestinal gas determined by means of an intestinal washout technique. \textit{N Engl J Med}. 1971;284(25):1394--1398.
\bibitem{hasler2006}
Hasler WL. Irritable bowel syndrome and bloating. \textit{Gastroenterol Clin North Am}. 2006;35(2):307--317.
\bibitem{azpiroz2005}
Azpiroz F, Malagelada JR. The pathogenesis of bloating and abdominal distension. \textit{Gastroenterology}. 2005;129(3):1060--1071.
\bibitem{mcmahon2014}
McMahon M, O'Morain CA. Flatulence and gaseousness. In: \textit{Evidence-Based Gastroenterology}. Elsevier; 2014.
\bibitem{winchell1995}
Winchell MJ, Robbins RC. Flatulence. In: \textit{Misoprostol in Gastroenterology}. 1995.
\bibitem{suarez1997}
Suarez F, Levitt MD. Intestinal gas. \textit{Curr Opin Gastroenterol}. 1997;13(2):130--134.
\end{thebibliography}
\clearpage
